\documentclass{article}
\usepackage{multicol}
\usepackage[a4paper]{geometry}
\usepackage{fancyhdr}
\pagestyle{fancy}
\lhead{Freihandel und Protektionismus}
\rhead{März 2026}
\begin{document}
\section{Freihandel und Protektionismus}
Die Handelsstrategie einer Nation kann in den \emph{Freihandel} und den \emph{Protektionismus} aufgeteilt werden.
 
\raggedcolumns
\begin{multicols}{2} 
\subsection{Freihandel}
Beim Freihandel ist, die der Name sagt, der Handel möglichst frei. So gibt es nur wenige bis garkeine Zölle, Qualitätsanforderungen oder andere Handelsvorschriften. Dies führt möglichst effizient zu einem möglichst hohen Wohlstand.
 
Dabei gibt es aber auch viele Probleme; der Freihandel beschleunigt die Umverteilung von unten nach oben, vereinfacht Lohndumping, hemmt neue Industrien eines Landes durch Kompetition aus dem Ausland und ist für die Sicherheit, werden sich z.\,B. Rüstungskonzerne angeguckt, verheerend und schafft nationale Abhängigkeiten.
 
\columnbreak 
 
\subsection{Protektionismus}
Beim Protektionismus hingegen geht es, wie der Name bereits sagt, darum, die Wirtschaft der eigenen Nation zu beschützen. Dies kann passieren indem
\begin{list}{}{\leftmargin=0cm}
 \item überbillige Produkte aus dem Ausland durch Handelshemmnisse gehemmt werden \emph{oder}
 \item Ware aus dem Inland gefördert wird, durch Subventionen oder staatliche Aufträge.
\end{list}
So wird der Wettbewerb gehemmt und der letztendliche Gesamtwohlstand ist geringer, die anderen genannten Probleme kommen aber auch nur in verringerter Form vor. 
\end{multicols}
 
\subsection{Handelshemmnisse}
\emph{Handelshemnisse} können in zwei Arten aufgeteilt werden; \emph{tarifäre} und \emph{nichttarifäre} Handelshemnisse.
\begin{description}
 \item[tarifäre] Handelshemnisse beinhalten alle Handelshemnisse, welche direkte finanzielle Abgaben beim Import darstellen, also alle Arten von Zöllen.
 \item[nichttarifäre] Handelshemnisse sind alle anderen arten an Handelshemnissen, beispielsweise Mengenbegrenzungen oder Qualitätsanforderungen, wie Normen und Standards.
\end{description}
Die Europärische Union hat durchaus nichttarifäre Handelshemnisse, wobei in der letzten Zeit, vorallem durch die zurzeitische Amerikanische Regierung, tarifäre relevanter wurden.
 
\subsection{Die Europäische Union hier}
Die EU hat eine gemeinsame Handelspolitik, welche den Handel und den Tausch von waren mit anderen Ländern regelt. So handelt die EU auch in Bezug auf Zölle als Union.
 
Innerhalb der Zollunion der EU ist der Handelsverkehr frei.
 
Die Europäische Kommission hat das Recht, Handelsabkommen mit anderen Ländern abzuschließen. 
\end{document}